\section{方法}
\label{sec:method}

\subsection{什么算捷径}
\label{sec:def}

如果一条只使用智能体可得信息、并且不执行基准所描述之调查的固定规则能给出正确答案，就说这道题存在\emph{捷径}。定义包含三点。规则必须是\emph{固定的}：写一次，应用于所有题目，不按题调整。规则只能使用\emph{智能体看得到的东西}：提示词，以及智能体能查询的数据；读取答案或生成图的规则只能给出上界。规则必须\emph{跳过设想的工作}：对 ExCyTIn 是从起点告警到终点告警的多跳跳转；对 SIABench 是分析抓包；对 GUIDE 是利用事件本身的证据；对 OTRF 是把事件与其上下文关联起来。

按信息来源，我们区分四类捷径。
\begin{itemize}\setlength\itemsep{1pt}
\item[\textbf{S1}] \emph{目标泄漏}：题目点出了答案所在的证据。生成的题目为了不产生歧义，会描述生成路径的终点，而这段描述足以找到它。
\item[\textbf{S2}] \emph{来源决定标签}：标签由谁打的决定，而不是由事件本身决定，例如由组织的定级政策决定。
\item[\textbf{S3}] \emph{生成痕迹}：数据的产生方式留下了区分类别的线索，例如攻击网络的地址或攻击工具的默认名称，而良性数据来自另一个更容易的分布。
\item[\textbf{S4}] \emph{构造的确定性}：在基准作者构建的环境中，对证据模型执行固定程序就能解出所有案例，基准因而无法区分会调查的模型和一段脚本。
\end{itemize}
这几类并不互斥，区别在于谁能修复：S1 和 S4 可以通过改变题目或环境的生成方式修复；S2 和 S3 需要新的数据。

\subsection{不调查的基线}
\label{sec:baselines}

对每个基准，我们按捷径允许的最省事方式写一条规则，不含 LLM。

\paragraph{ExCyTIn} 两条规则都通过读取题目点名的告警来作答。\emph{图基线}按题目与告警名的词重叠（去掉停用词）对事件图中的告警排序，取第一名，回答它的第一个类型与题目所问类型相符的相邻实体（固定关键词表把"IP address"映射到 IP 实体，"account"映射到账户实体，等等），并排除题目中已出现的值。\emph{告警表基线}做同样的事，但只使用 \texttt{SecurityAlert} 表的 \texttt{AlertName} 和 \texttt{Entities} 两列，任何智能体都能查询这张表。智能体看不到图；图基线衡量点名线索能解释基准的多少，告警表基线表明智能体能够利用它。两者都只给一个答案，规范化后按字符串精确匹配计分。ExCyTIn 本身用 LLM 判分并给中间步骤部分分；我们的计分更严，所以比较对模型有利。我们还为每道题记录两个属性：终点告警是否按词重叠排第一（\emph{点名}），以及标准答案是否为终点告警的实体（\emph{相邻}）。

\paragraph{SIABench} 规则：当且仅当告警文字涉及 172.16.0.1 时回答"真阳性"。172.16.0.1 是 CIC-IDS2017 中攻击网络进入受害网络的网关~\citep{cicids2017}。同时报告多数类。

\paragraph{GUIDE} 规则：回答同一组织此前对同一检测器最常给出的定级。我们在训练集上按组织内的时间切分测量，并与只用事件自身元数据、按组织不交叉切分训练的分类器比较。

\paragraph{OTRF} 在 27 个远程执行录制中，把每次服务安装和计划任务的创建或修改视为一条告警，并按录制元数据打标签。规则：任务位于 \texttt{\textbackslash Microsoft\textbackslash Windows\textbackslash} 下或为 OneDrive 更新任务，或服务程序为 \texttt{svchost -k} 或 Defender 组件，则判为良性。

\paragraph{构建的环境} 我们还把方法用于此前为 LLM 分诊智能体研究构建的三个模拟分诊环境：\emph{sec-triage}（72 个网络与主机告警案例，良性与恶意原因相互竞争）、\emph{sigma-triage}（由 SigmaHQ 检测规则派生的 152 个案例~\citep{sigmahq}）和 \emph{comp-triage}（72 个案例，每个攻击都与使用相同工具的合法活动配对）。不调查的基线是不含 LLM 的控制器：按期望信息增益选择探针，当一条观测明确确认领先原因时停止。

\subsection{流程}
\label{sec:protocol}

ExCyTIn 的规则在基准最新版本的 418 道训练题上写定并冻结，同时写下假设，然后在其 599 道测试题上只运行一次（\cref{app:protocol}）。之后我们发现 ExCyTIn 论文的成绩来自更早的版本，该版本由相同事件生成、有 589 道测试题，而且作者公开了 14 次模型运行在该版本上的逐题日志。在读取任何日志之前，我们登记了补充方案：在早期版本上运行同样冻结的规则，并比较每个模型在有捷径和无捷径题目上的成功率。SIABench 的规则是看过数据之后发现的，作为事后观察报告；它是对数据集一项有文档记载之性质的单次比较。GUIDE 只用训练集，OTRF 只用我们下载的 27 个录制；GUIDE 测试集没有读取。

\paragraph{统计} ExCyTIn 只有八个事件，同一事件的题目共享告警和实体。因此我们按事件而不是按题目重抽：95\% 区间来自 2{,}000 次按事件聚类的 bootstrap（种子 2040）。跨模型使用单侧符号检验。不是我们自己计算的模型成绩均抄自基准论文，并标注为报告值。
